\begin{explanationbox}
	\textbf{\textcolor{CExplanation}{一句话直觉}}
	双曲线中 $c$ 相当于直角三角形斜边，$a$ 和 $b$ 分别是两条直角边。

	\medskip
	\textbf{\textcolor{CExplanation}{核心拆解}}
	\begin{description}[style=nextline, leftmargin=2.8em, labelsep=0.5em, itemsep=0.45em]
		\item[\textbf{要点一（定义导出关系）：}]
			由双曲线定义 $|PF_1-PF_2|=2a$，通过坐标运算化简可得 $c^2=a^2+b^2$。
		\item[\textbf{要点二（与椭圆对比）：}]
			椭圆中 $c^2=a^2-b^2$（$a$ 最大），双曲线中 $c^2=a^2+b^2$（$c$ 最大）。
	\end{description}

	\medskip
	\textbf{\textcolor{CExplanation}{几何本质（直角三角形）}}
	实半轴 $a$、虚半轴 $b$ 和半焦距 $c$ 构成一个直角三角形，$c$ 是斜边。

	\medskip
	\textbf{\textcolor{CExplanation}{代数意义（公式推导结论）}}
	$c^2=a^2+b^2$ 是从双曲线标准方程推导中直接得到的代数关系，也是连接 $a,b,c$ 的桥梁。

	\medskip
	\textbf{\textcolor{CExplanation}{考点价值（常见应用）}}
	\begin{itemize}[leftmargin=*, itemsep=0.5em, topsep=0.4em]
		\item \textbf{考法一（求标准方程）：} 已知焦点坐标和实轴长，利用 $c^2=a^2+b^2$ 求解 $b$。
		\item \textbf{考法二（求离心率）：} 直接由 $e=\frac{c}{a}=\sqrt{1+\left(\frac{b}{a}\right)^2}$ 计算。
	\end{itemize}

	\medskip
	\textbf{\textcolor{CExplanation}{顿悟点}}
	双曲线中 $c$ 最大，椭圆中 $a$ 最大，这是区分二者的关键几何特征。

	\medskip
	\textbf{\textcolor{CExplanation}{使用场景}}
	\begin{itemize}[leftmargin=*, itemsep=0.5em, topsep=0.4em]
		\item \textbf{场景一（直接求参数）：} 已知标准方程，直接读取 $a,b$ 并计算 $c$。
		\item \textbf{场景二（逆向设方程）：} 已知焦点和实轴长，设标准方程并利用关系求 $b$。
	\end{itemize}

\end{explanationbox}
